\subsection{The symplectic group}

Recall that we have defined a groupoid whose objects are symplectic spaces, and whose isomorphisms are the symplectomorphisms.
And \refmath[corollary]{dimsympinvar} says that the dimension of a symplectic space is its only significant invariant.
Let us denote by
$$ \Sp(V,\omega) = \set{\Psi\in\gl(V)}[\Psi^*\omega=\omega] $$
the automorphism group of $(V,\omega)$, called the {\emphcolor symplectic group}.

By \refmath[proposition]{eucsplie}, $\Sp(V,\omega)$ is a Lie group.
Following this proposition, we have the following result.

\bthrm

    The symplectic group $\Sp(V,\omega)$ is a Lie group with Lie algebra
    $$ \lsp(V,\omega) = \set{T\in\endo(V)}[\omega(A\blank,\blank)+\omega(\blank,A\blank)=0] . $$

\ethrm

We can understand this group by focusing on the case of $(\bR^{2n},\omega_0)$.
Then $\Sp(2n)=\Sp(\bR^{2n},\omega_0)$ may be thought of as a Lie group of matrices $\Psi\in\Mat_{2n}$ satisfying $\Psi^\top J_0\Psi=J_0$.
Such matrices are called {\emphcolor symplectic}.

\blemm[name={sympdet1}]

    Symplectic matrices have determinant one.

\elemm

\bproof

    By \refmath[corollary]{ndextpower}, $\omega_0^n$ is nonzero and thus a volume form on $\bR^{2n}$.
    Furthermore a symplectic matrix $\Psi$ satisfies $\Psi^*\omega_0=\omega_0$ and thus $\Psi^*\omega_0^n=\omega_0^n$.
    Because this is a volume form, $\det\Psi=1$.
    \qed

\eproof

We identify $\bR^{2n}$ with $\bC^n$ by identifying $(x,y)\in\bR^{2n}$ with $x+iy\in\bC^n$.
Then $J_0$ corresponds to multiplication by $i$ in $\bC^n$.
With this identification, $\gl_n(\bC)$ is the subgroup of $\gl_{2n}(\bR)$ consisting of matrices which commute with $J_0$:
$$ \gl_n(\bC) = \set{A\in\gl_{2n}(\bR)}[J_0A=AJ_0] . $$
In this identification, a matrix $X+iY\in\Mat_n(\bC)$ corresponds to the matrix
$$ X + iY\in\Mat_n(\bC) \oto \pmatrix{X&-Y\cr Y&X} \in \Mat_{2n})(\bR) . $$
Notice how the conjugate transpose acts on $\Mat_n(\bC)$:
$$ (X+iY)^\dagger = (X-iY)^\top = \pmatrix{X^\top &Y^\top\cr -Y^\top&X^\top} = \pmatrix{X&-Y\cr Y&X}^\top . $$
Thus $X+iY\in\unitary(n)$ iff
$$ \pmatrix{X&-Y\cr Y&X}^{-1} = \pmatrix{X^\top &Y^\top\cr -Y^\top&X^\top} , $$
which is iff
$$ XX^\top + YY^\top = I,\qquad XY^\top = YX^\top . $$

Notice that
$$ i(X+iY)^\dagger(X+iY) = i \iff (X+iY)^\dagger = (X+iY)^{-1} \iff X+iY\in\unitary(n) . $$
Thus we see
$$ \unitary(n) = \gl_n(\bC)\cap\Sp(2n) . $$
In fact we can extend this result.

\blemm

    $$ \Sp(2n)\cap\orth(2n) = \Sp(2n)\cap\gl_n(\bC) = \orth(2n)\cap\gl_n(\bC) = \unitary(n) . $$

\elemm

\bproof

    A matrix $\Psi\in\gl_{2n}(\bR)$ satisfies
    $$ \Psi \in \gl_n(\bC) \iff \Psi J_0=J_0\Psi,\qquad \Psi\in\Sp(2n) \iff \Psi^\top J_0\Psi = J_0,\qquad \Psi\in\orth(2n) \iff \Psi^\top\Psi = I . $$
    Note that any of these two conditions imply the third, so we have the first two equalities.
    By the above discussion, we also get the final.
    \qed

\eproof

For a linear operator $\Psi$ we denote its spectrum by $\sigma(\Psi)$; this is its set of eigenvalues.

\blemm

    Let $\Psi\in\Sp(2n)$, then $\lambda\in\sigma(\Psi)$ iff $\lambda^{-1}\in\sigma(\Psi)$, and the multiplicities of $\lambda$ and $\lambda^{-1}$ are equal.
    If $\pm1$ is an eigenvalue of $\Psi$ then it has even multiplicity.

    Finally, if $\lambda,\lambda'\in\sigma(\Psi)$ are two (not necessarily distinct) eigenvalues such that $\lambda\lambda'\neq1$, then $\omega_0(z,z')=0$
    for any eigenvectors of eigenvalues $\lambda,\lambda'$ respectively.

\elemm

\bproof

    The first statement is due to the fact that $\Psi^{-1}$ is similar to $\Psi^\top$:
    $$ \Psi^\top = J_0\Psi^{-1}J_0^{-1} , $$
    $\Psi$ and $\Psi^\top$ share the same eigenvalues (with the same multiplicities), and $\lambda$ is an eigenvalue of $\Psi$ iff $\lambda^{-1}$ is of $\Psi^{-1}$ (and the multiplicities agree).

    Thus the total multiplicity of the eigenvalues not equal to $\pm1$ is even (as they come in pairs).
    The determinant of $\Psi$ is the product of eigenvalues, as they come in pairs with their inverses this is equal to $(-1)^m$ where $m$ is $-1$'s multiplicity.
    By \refmath[lemma]{sympdet1} this is equal to $1$ so $m$ is even.
    Because $\Psi$ has dimension $2n$, the multiplicity of $1$ must be even as well.

    The last statement is because
    $$ \omega_0(z,z') = \omega_0(\Psi z,\Psi z') = \lambda\lambda'\omega_0(z,z') $$
    and since $\lambda\lambda'\neq1$, $\omega_0(z,z')=0$.
    \qed

\eproof

Suppose $D$ is a diagonal matrix, $D=\diag(\lambda_1,\dots,\lambda_n)$, then define
$$ D^\alpha = \diag(\lambda_1^\alpha,\dots,\lambda_n^\alpha) $$
for any real $\alpha\in\bR$ where this is defined.
Suppose $A\in\Mat_n(\bR)$ is a diagonalizable matrix with $A=P^{-1}DP$, $D$ diagonal.
Then define
$$ A^\alpha = P^{-1}D^\alpha P $$
for real $\alpha$ where this is defined.
Note that this agrees with $A^n$ for $n$ natural.
It is easy to see that this is well-defined (i.e.\ not dependent on $P,D$).

If $A^\alpha$ exists and $Av=\lambda v$, then $A^\alpha v=\lambda^\alpha v$.
This is clear for diagonal $A$, and extends to diagonalizable $A$ by extending $v$ to an eigenbasis and using that $P$.

Note that if $A$'s eigenvalues are all positive, then this is defined for all $\alpha$.
In particular if $S$ is symmetric and positive definite, then its powers exist for all real $\alpha$.

\blemm[name={expiscont}]

    Exponentiation is continuous as a map
    $$ \Sym_+\times\bR\to\Sym_+,\qquad (S,\alpha)\mapsto S^\alpha $$
    where $\Sym_+\subseteq\Mat_n$ is the subspace of positive-definite symmetric matrices.

\elemm

\bproof

    Let $S\in\Sym_+$, and let its minimum and maximum eigenvalues by $\lambda_{\min},\lambda_{\max}$ respectively.
    Then we know for all vectors $v$,
    $$ \lambda_{\min}\norm v^2 \leq v^\top Sv\leq \lambda_{\max}\norm v^2 $$
    and these bounds are tight.
    We claim that there is an open neighborhood $S\in U\subseteq\Sym_+$ and numbers $0<m<M$ such that every $A\in U$ has eigenvalues bound by $m,M$.

    Let $m=\lambda_{\min}/2$, then for every symmetric matrix $A$ and $v$ of unit norm:
    $$ v^\top Av = v^\top Sv + v^\top(A-S)v \geq \lambda_{\min} - \abs{v^\top(A-S)v} \geq \lambda_{\min} - \norm{(A-S)v} \geq \lambda_{\min} - \norm{A-S} $$
    where $\norm{A-S}$ denotes the operator norm.
    Thus if $\norm{A-S}<m$ then the right-hand side is equal to $\lambda_{\min}/2$ and thus $v^\top Av\geq\lambda_{\min}/2$; so all of $A$'s eigenvalues are bound
    from below by $m$.
    So all matrices in the ball of radius $m$ around $S$ have eigenvalues bound from below by $m$.

    We can bound this from above too: for unit norm $v$
    $$ v^\top Av \leq \norm A \leq \norm S+\norm{A-S} , $$
    and so if we require $\norm{A-S}<1$ then
    $$ v^\top Av \leq \norm S + 1 = \lambda_{\max} + 1 , $$
    so matrices in the ball of radius $1$ around $S$ have eigenvalues bound from above by $M=\lambda_{\max}+1$.
    So the desired open set is
    $$ U = \set{A\in\Sym_+}[\norm{A-S}<1,m] . $$

    Instead of considering exponentiating over arbitrary $\alpha$, we show it is continuous for $a\leq\alpha\leq b$ for any arbitrary $a,b$; this is clearly sufficient.
    Consider the map $f\colon[m,M]\times[a,b]\to\bR$ by $f(x,t)=t^x$, this is continuous as $m>0$.
    By Stone-Weirestrass, because the codomain is compact, we can approximate $f$ uniformly by a sequence of polynomials $p_k\in\bR[x,t]$.

    Now suppose $A\in U$, so $A=P^{-1}\diag(\lambda_1,\dots,\lambda_n)P$, and because $p_k$ are polynomial
    $$ p_k(x,A) = P^{-1}\diag(p_k(x,\lambda_1),\dots,p_k(x,\lambda_n))P . $$
    Thus
    $$ A^\alpha - p_k(\alpha,A) = P^{-1}\diag(\dots,\lambda_i-p_k(\alpha,\lambda_i),\dots)P . $$
    Because $A\in U$ the $\lambda_i$ fall in $[m,M]$, and so under the operator norm
    $$ \norm{A^\alpha-p_k(\alpha,A)} = \max_i\abs{\lambda_i-p_k(\alpha,\lambda_i)} \leq \max_{t\in[m,M]}\abs{t^\alpha-p_k(\alpha,t)} . $$

    And so we see that the map $F\colon U\times[a,b]\to\Sym_+$, $F(A,\alpha)=A^\alpha$, satisfies for all $A,\alpha$
    $$ \norm{F(A,\alpha) - p_k(\alpha,A)} \leq \max_{t\in[m,M],\alpha\in[a,b]}\abs{t^\alpha-p_k(\alpha,t)} , $$
    so that $(\alpha,A)\mapsto p_k(\alpha,A)$ uniformly converges to $F$.
    Because $p_k$ are all polynomial and thus continuous, so too is $F$.
    \qed

\eproof

\blemm

    If $S\in\Sp(2n)$ is symmetric and positive definite, then $P^\alpha\in\Sp(2n)$ for all $\alpha\in\bR$.

\elemm

\bproof

    As $S$ is symmetric and positive definite we can decompose $\bR^{2n}$ into eigenspaces
    $$ \bR^{2n} = \bigoplus_{i=1}^k E_i,\qquad E_i = \ker(\lambda_iI-S) , $$
    for $\lambda_i>0$ the eigenvalues.
    For $z,w\in\bR^{2n}$ let their projections into $E_i$ be denoted $z_i,w_i$ respectively.
    Then because $S$ is symplectic
    $$ \omega_0(z_i,w_j) = \omega_0(Pz_i,Pw_j) = \lambda_i\lambda_j\omega_0(z_i,w_j) . $$
    So either $\lambda_i\lambda_j=1$ or $\omega_0(z_i,w_j)=0$, and thus
    Then
    $$ \omega_0(P^\alpha z,P^\alpha w) = \sum_{i,j}^k(\lambda_i\lambda_j)^\alpha\omega_0(z_i,w_j) = \sum_{i,j}^k\omega_0(z_i,w_j) = \omega_0(z,w) $$
    so that $P^\alpha\in\Sp(2n)$.
    \qed

\eproof

\bprop

    \benum
        \item $\unitary(n)$ is a maximal compact subgroup of $\Sp(2n)$.
        \item The inclusion $\unitary(n)\to\Sp(2n)$ is a strong deformation retract.
        In particular $\Sp(2n)$ is connected.
    \eenum

\eprop

\bproof

    We first prove (2).
    Define
    $$ H\colon[0,1]\times\Sp(2n)\to\Sp(2n),\qquad H(t,\Psi) = \Psi(\Psi^\top\Psi)^{-t/2} . $$
    This is continuous by \refmath[lemma]{expiscont}, as $\Psi^\top\Psi$ is positive definite for any non-singular $\Psi$.
    This is well-defined by the previous lemma, $H(\Psi,t)\in\Sp(2n)$ for all $t$.

    Now
    $$ H_0 = \id,\qquad H_t|_{\unitary(n)} = \id $$
    where the latter is because $\Psi^\top\Psi=I$ for unitary $\Psi$.
    And for $\Psi\in\Sp(2n)$, $H_1(\Psi)$ is orthogonal:
    $$ H_1(\Psi)^\top H_1(\Psi) = (\Psi^\top\Psi)^{-1/2}\Psi^\top\Psi(\Psi^\top\Psi)^{-1/2} = I . $$
    Because $\Sp(2n)\cap\orth(2n)=\unitary(n)$, we have $H_1(\Sp(2n))=\unitary(n)$.
    Thus $H$ is indeed a strong deformation retract.
    And so $\Sp(2n)$ and $\unitary(n)$ are homotopy equivalent, and $\unitary(n)$ is connected so therefore $\Sp(2n)$ is as well.

    We show (1) now.
    Suppose $G$ is a compact subgroup of $\Sp(2n)$ containing $\unitary(n)$, but not equal to $\unitary(n)$.
    Then let $\Psi\in G-\unitary(n)$, and so $H_1(\Psi)=\Psi(\Psi^\top\Psi)^{-1/2}\in\unitary(n)\subseteq G$.
    Thus the positive definite symmetric matrix $P=\Psi^\top\Psi$ belongs to $G-\unitary(n)$.
    Therefore it must have an eigenvalue greater than $1$: its eigenvalues come in pairs $\lambda,\lambda^{-1}$, so otherwise $P=I$.
    But then the sequence $P,P^2,P^3,\dots$ in $G$ has no convergent subsequence, contradicting $G$'s compactness.
    \qed

\eproof

\bremark

    Every symplectic matrix $\Psi\in\Sp(2n)$ can be uniquely written as $\Psi=UP$ where $U\in\unitary(n)$ and $P$ is a symmetric, positive definite, symplectic matrix.
    Existence is by
    $$ U = \Psi(\Psi^\top\Psi)^{-1/2},\qquad P = (\Psi^\top\Psi)^{1/2} . $$
    Uniqueness is because the only symmetric positive definite unitary matrix is $I$ (because its eigenvalues come in pairs, and thus must all be $1$).
    If $U\in\unitary(n)$ can be written as $U=U'P'$, then $P=U(U')^{-1}$ is unitary as well, and thus $P=I$ so $U=U'$.
    So this is true for unitary matrices, and thus if $UP=U'P'$ we have $U=U'P'P^{-1}$ so $U=U'$ and $P'P^{-1}=I$ meaning $P=P'$.
    This shows uniqueness.

    This decomposition is called {\emphcolor symplectic polar decomposition}.

\eremark

\bprop

    The complex determinant map $\det_\bC\colon\unitary(n)\to S^1$ induces an isomorphism of fundamental groups $\pi_1(\unitary n)\cong\pi_1(S^1)\cong\bZ$.

\eprop

\bproof

    The determinant map is a fibration with typical fiber $\sunitary(n)$.
    Indeed for every $e^{i\theta}\in S^1$ we can consider the open subset $W$ of $S^1$ consisting of numbers whose arguments are between $\theta-\epsilon$ and $\theta+\epsilon$.
    Then $\det_\bC^{-1}(W)\to\sunitary(n)\times W$ given by $\Psi\mapsto(e^{i\phi/n}\Psi,e^{i\phi})$ where $e^{i\phi}=\det\Psi$ is a homeomorphism with inverse $(\Phi,e^{i\phi})\to e^{i\phi/n}\Phi$.
    It also commutes with $\det_\bC$ because $\det(e^{i\phi/n}\Phi)=e^{i\phi}$.

    Thus there is an exact sequence
    $$ \pi_1(\sunitary(n)) \to \pi_1(\unitary n) \to \pi_1(S^1) \to \pi_0(\sunitary(n)) . $$
    It is sufficient then to show that $\sunitary(n)$ is simply connected, and we get the desired result.

    It is trivial that $\sunitary(n)$ is simply connected for $n=1$, as $\sunitary(1)=\set1$.
    For $n\geq2$ the map $\sunitary(n)\to S^{2n-1}$ which maps $U\in\sunitary(n)$ to its first column is a fibration with typical fiber $\sunitary(n-1)$.
    And so we have an exact sequence
    $$ \pi_1(\sunitary(n-1)) \to \pi_1(\sunitary n) \to \pi_1(S^{2n-1}) \to \pi_0(\sunitary(n-1)) . $$
    By induction this means $\pi_1(\sunitary n)\cong\pi_1(S^{2n-1})=1$ because $2n-1>1$.
    \qed

\eproof

